%%%%%%%%%%%%%%%%%%%%%%%%%%%%%%%%%%%%%%%%%%%%%%%%%%%%%%%%%%%%%%%%%%%%%
%  sub/en-1-introduction.tex  —  1. Introduction
%%%%%%%%%%%%%%%%%%%%%%%%%%%%%%%%%%%%%%%%%%%%%%%%%%%%%%%%%%%%%%%%%%%%%
\section{Introduction}\label{sec:introduction}

% 서론에서는 연구가 다루는 과학교육 문제와 그 중요성을 제시하고, 선행연구를
% 비판적으로 검토하여 연구 공백(gap)을 드러낸 뒤, 본 연구의 목적과 연구 질문을
% 명확히 밝힙니다. Research in Science Education 은 이론적 기여와 실증적 근거의
% 연결을 중시하므로, 연구 질문이 어떤 이론적·실천적 함의를 겨냥하는지 초반에
% 분명히 하세요.

Write the introduction here. State the science-education problem, review prior
work to expose the gap, and end with explicit research questions.

% 인용은 \parencite{key} (괄호 인용) 또는 \textcite{key} (본문 인용)로 씁니다.
% 예: Prior work has examined this issue \parencite{example2024}, and
% \textcite{example2024} argued that \dots
